\section{Dodecaphasic record and rational phase constant}
\label{sec:k-registro}

The generation of regional coordinates does not exhaust the information in the
APP--TRIT--TPK state. One Archimedean reading determines a value through
compatible cylinders; another reading preserves the opposition of the sheets, the
orientation of the route, and the record of the incidences traversed. The
dodecaphasic record belongs to this second reading before entering the
closure of alpha. It is therefore not defined by subtracting a physical constant from
a combination of previously known values. It is reconstructed from a
signed coordinate of the same arithmetic state with trajectory memory that produces the regional
coordinates.

The distinction is necessary even when two objects have twelve
coordinates. A twelve-block word from the catalog, a list of twelve
oriented events, a signed integer vector, and a base-thousand dodecaphasic record are
different objects. Comparing them requires the maps that
relate them. Equality of length does not provide those maps, and
the loss of signs is not a mere simplification of notation.

This chapter maintains the causal cut fixed in the preceding chapters:
the nine APP positions on the horizontal and vertical axes produce the
arithmetic sheets; TRIT types regime and orientation; TPK selects,
transports, updates, and prolongs, preserving remainder, quotient, carry,
sheet, route, boundary, and memory. The coordinates used here are obtained
from the joint discrete structure of the continuum. The chain
\(w_6\to w_{12}\to w_{18}\to w_{24}\to w_{30}\to R_{36}\to G_9\)
is not replaced by a finite calendar. Nor is the order
\(U_t\to\Gamma_9\to K_{\mathrm{ph}}\) reversed: the last map is a
reduced-memory reading, not the generator of the signed record.

% Proposed addition at the beginning of registro_k.tex, after its introduction.
% Preserves all current definitions, formulas, and proofs.
% Provenance: registro_k.tex, §§k-hadamard/k-transversal/k-kappas;
% k_reversibilidad.tex; incidencia_articulacion.tex, inc-ampl-cuaterna.
\paragraph{Reversible coordinatization and incidence determination.}
\label{ampl-alfa-k-intro-reversible}
The result organizing this chapter is the equivalence between three
coordinatizations of a record previously produced by
APP--TRIT--TPK. The Hadamard transformation relates the signed record
to the integer record; the difference extractor relates
the latter to its transverse variations and its uniform component:
\[
 U\underset{\mathcal H_{12}}{\overset{\mathcal H_{12}/4}{\longleftrightarrow}}K
 \overset{(D_3,D_4,Q)}{\longleftrightarrow}
 (D_3K,D_4K,Q(K)).
\]
The first equivalence is established on the integral sublattice
$\mathcal L_H$; the second, on the image of the extractor. With
base $1000$, the $12$ positions, and the phase origin fixed, the evaluation
$K\mapsto\kappa_{\mathrm{per}}$ adds a fourth reversible
coordinatization: multiplication by $1000^{12}-1$ followed by Euclidean
division recovers the complete record. The proofs in
\S\ref{subsec:k-hadamard}, \S\ref{subsec:k-transversal}, and
\S\ref{subsec:k-kappas} establish these equivalences in their respective
domains.

Two compositions subsequently act on this record. The positional
composition combines its components with the closure, propagation,
and self-scaling coordinates, and its normalization determines $\alpha_A$.
The incidence composition applies the cubic threshold and yields
\[
 K\longmapsto B_K=\{j:K_j\geq729\}=\{4,6,9,10\}.
\]
The complete record allows both compositions to be reconstructed. The support
$B_K$ specifically preserves membership in the class selected
by the threshold. The identity $B_K=H_e\cap P_\pi$ in
\eqref{inc-ampl-cuaterna} links this integer extraction to the incidence
of the regional codes. The negative support of their balance with $K$
completes the flag $B_K\subset H^-_\alpha$ developed in the
exceptional extension. Thus, the function of $K$ is defined through
an invertible transformation and precise subsequent maps;
the additional information on depth and transport remains
in the state from which the record originates.



\subsection{Cyclic extension and memory grading}
\label{subsec:k-calendario}

The observable calendar contains one hundred eight positions, organized into
twelve nonadic windows. With positive indices, let
\[
 E_m=\{9(m-1)+1,\ldots,9m\},\qquad 1\leq m\leq12.
\]
The ordered support \(E_1\sqcup\cdots\sqcup E_{12}\) preserves position and
window membership. The group \(C_{108}=\mathbb Z/108\mathbb Z\), by contrast,
preserves the law of cyclic translation. For the representative
\(0\leq\widetilde\theta<108\), the coordinates
\[
 \beta(\theta)=\left[\left\lfloor\frac{\widetilde\theta}{9}\right\rfloor\right]_{12},
 \qquad r(\theta)=1+(\widetilde\theta\bmod9)
\]
separate the dodecaphasic sector and the interior nonadic position. Neither
by itself preserves the total number of turns.

\begin{proposition}[Calendar carry]
\label{prop:k-extension}
The maps
\[
 \iota:C_{12}\longrightarrow C_{108},\quad[b]\longmapsto[9b],
 \qquad p:C_{108}\longrightarrow C_9,\quad[\theta]\longmapsto[\theta]_9
\]
form a nonsplit exact sequence
\[
 0\longrightarrow C_{12}\xrightarrow{\iota}C_{108}
 \xrightarrow{p}C_9\longrightarrow0.
\]
For the set-theoretic section choosing the representatives
\(0,\ldots,8\), the defect of additivity is the carry
\[
 c(r,r')=\left[\left\lfloor\frac{\widetilde r+\widetilde r'}9\right\rfloor\right]_{12}.
\]
\end{proposition}

\begin{proof}
The kernel of \(p\) consists of the multiples of nine and equals
the image of \(\iota\). Reduction is surjective, and multiplication
by nine is injective on the indicated domain. If the sequence
split, the middle group would be isomorphic to \(C_{12}\times C_9\), whose
exponent is \(\operatorname{lcm}(12,9)=36\); the exponent of
\(C_{108}\) is 108. This contradiction proves nonsplitting. Finally,
Euclidean division of \(\widetilde r+\widetilde r'\) by nine gives the
displayed carry and verifies the section identity.
\end{proof}

The result explains a difference that will recur in alpha: returning to the
phase is not equivalent to returning to the state. If \(q_{\mathrm{mem}}\) counts
nonadic turns, the calendar retains only
\(\beta=[q_{\mathrm{mem}}]_{12}\). After one hundred eight steps,
\(\beta\) returns, but \(q_{\mathrm{mem}}\) increases by twelve. The
record also preserves the cylinder, the incidence ledger, and vertical
memory. Closing a finite chart does not turn that history into a
memoryless cycle.

\subsection{The oriented record and the integral \texorpdfstring{\(1+5+6\)}{1+5+6} chart}
\label{subsec:k-registro-integral}

For a history produced by TPK, let \(\tau_m\) be the subroute over
\(E_m\), \(\sigma_m\) its orientation, \(\kappa_m\) the boundary
carry, and \(\Lambda_m\) the local incidence ledger. The record is
\[
 \mathcal R_{12}(x)=
 \bigl((E_m,\tau_m,\sigma_m,\kappa_m,\Lambda_m)\bigr)_{m=1}^{12}.
\]
Its domain is the space of arithmetic histories with trajectory memory, not the set of
phase labels. Projection onto the calendar forgets precisely some
of the data required by the signed reader.

An initial chart of the record can be written through oriented
events. The digital reader fixes the set of event codes
\(\{2,4,6,8\}\) and the sector orientation
\[
 \Sigma_{12}=(+,+,+,-,-,-,+,+,+,-,-,-).
\]
If \(d_t\) is the symbol of the digital record at position \(t\), one obtains
\[
 e_i=\Sigma_{12}(i+1)
 \#\{t\in E_{i+1}:d_t\in\{2,4,6,8\}\},
 \qquad0\leq i<12.
\]
The formula keeps the nonnegative event count separate from
orientation. The digital symbol, cardinal direction, and TRIT sign
retain their respective domains. The
vector \(e\), whose components are signed local counts, is not yet
the vector \(U\) appearing below either.

Opposite positions determine, for \(0\leq i<6\),
\[
 p_i=e_i+e_{i+6},\qquad a_i=e_i-e_{i+6}.
\]
With
\[
 s_C=\sum_{i=0}^{5}p_i,\qquad M_i=p_i-p_5\quad(0\leq i<5),
\]
the integral chart is defined by
\[
 \mathcal L_{12}(e)
 =(s_C,M_0,\ldots,M_4,a_0,\ldots,a_5).
\]
The \(1+5+6\) decomposition is not an arbitrary allocation of coordinates:
one coordinate preserves charge, five compare pairs relative to a
reference pair, and six preserve the opposition of the semicoronas.

\begin{lemma}[Inversion of the integral chart]
\label{lem:k-carta-integral}
An integer tuple \((s_C,M_0,\ldots,M_4,a_0,\ldots,a_5)\) belongs to the
image of \(\mathcal L_{12}\) if and only if
\[
 s_C-\sum_{i=0}^{4}M_i\equiv0\pmod6,
 \qquad p_i\equiv a_i\pmod2\quad(0\leq i<6),
\]
where
\[
 p_5=\frac{s_C-\sum_{i=0}^{4}M_i}{6},\qquad p_i=p_5+M_i\quad(i<5).
\]
In that case its preimage is unique and is given by
\[
 e_i=\frac{p_i+a_i}{2},\qquad e_{i+6}=\frac{p_i-a_i}{2}.
\]
\end{lemma}

\begin{proof}
The sum of the five margins is \(s_C-6p_5\), yielding
division by six. Once the \(p_i\) are reconstructed, each pair of equations
\(p_i=e_i+e_{i+6}\), \(a_i=e_i-e_{i+6}\) has the indicated solution.
This solution is integral exactly under the parity congruence. All
operations are inverted, so no residual freedom remains.
\end{proof}

This reversibility shows what information is preserved; it does not authorize
removing the rest of the ledger. The full record has a larger domain
than this count chart. Extracting the channels feeding
\(U\) uses that full record and does not follow from the list
\((e_i)\) through an identification of names.

% Working supplement. Does not replace registro_k.tex or close E108 by decree.
% Provenance and scope: INFORME_RECUPERACION.md, same folder.
\subsection{Incidence evaluation preceding the dodecaphasic record}
\label{subsec:registro-evaluacion-incidencial}

Recovering a record from its cyclic differences and
total charge is distinct from evaluating those coordinates
on the incidence ledger. The first problem is solved through the
transverse extractor of Proposition~\ref{prop:k-transversal}.
For the second, a historical presentation of the
record provides three integer counts in each window. It is useful
to make their composition explicit and determine exactly the information
their evaluation requires.

Let \(\mathscr L\) be a previously generated ledger of visits and oriented
traces. Each visit preserves the route, APP position, sheet,
instant, multiplicity, and boundary incidence. Arithmetic evaluation
produces the integer and its decomposition
\[
 n_t=r_t+9q_t,\qquad 1\leq r_t\leq9.
\]
The integer \(n_t\) is computed from the sum or product sheet; its
remainder is not received as an independent value. The distinction between
corona incidence and interior incidence remains explicit for visits
with remainder \(9\).

In window \(E_m\), the count presentation considers
\begin{align}
 A_m&=\sum_{v\in E_m}w(v)
             \mathbf1_{\{1,3,5,7\}}(r(v)),\\
 C_m&=\sum_{j\geq0}\bigl(\nu^-_{120,m}(j)-\nu^+_{120,m}(j)\bigr),\\
 V_m&=\sum_{v\in E_m}w(v)\mathbf1_{\{9\}}(r(v))\epsilon_\partial(v).
\end{align}
Here \(w(v)\) is the incidence multiplicity, and
\(\epsilon_\partial(v)=1\) on the corona, \(-1\) in the interior.
The traces counted by \(\nu^\pm\) must be enumerated through the
corresponding ledger rule. The integer sum requires finite support or an
alternative construction determining its meaning. The number of
chronological incidences and the reduced length \(120(1+3j)\) have
different types; their relation requires the reader's length unit.

Define \([x]_{10}\in\{0,\ldots,9\}\) as the Euclidean
representative and retain the quotients
\[
 x=[x]_{10}+10\lfloor x/10\rfloor.
\]
The incidence block is
\begin{equation}
 z_m=100[A_m]_{10}+10[C_m]_{10}+[V_m]_{10}.
 \label{eq:registro-bloque-incidencial}
\end{equation}
Composition yields the map
\begin{equation}
 \mathcal E_{\rm cnt}(\mathscr L)=
 \left((z_m-z_{m+3})_{m=1}^{12},
       (z_m-z_{m+4})_{m=1}^{12},\sum_{m=1}^{12}z_m\right).
\label{eq:registro-composicion-incidencial}
\end{equation}
Indices are read cyclically modulo \(12\). In this subsection,
\((D_jz)_m=z_m-z_{m+j}\), for \(j=3,4\), and
\(Q(z)=\sum_{m=1}^{12}z_m\); thus the preceding composition is
\((D_3z,D_4z,Q(z))\). This convention is retained in
\S\ref{subsec:k-transversal}, where complete integral reconstruction
is proved.
This formula produces its coordinates from the counts. It does not use
\(K\), \(U\), \(\alpha\), or the expected transverse coordinates
as arguments. Its identification with
\(\mathcal E_{108}^{90,120}\) requires composition with the effective collection and
incidences of the terminal state; it does not follow solely
from invertibility of the transverse system.

\begin{proposition}[Sufficient and necessary arithmetic information]
\label{prop:registro-36-residuos}
Let \(N,N'\in\mathbb Z^{12\times3}\), ordered by the counts
\((A,C,V)\). For the map defined by
\eqref{eq:registro-bloque-incidencial}--\eqref{eq:registro-composicion-incidencial},
\[
 \mathcal E_{\rm cnt}(N)=\mathcal E_{\rm cnt}(N')
 \quad\Longleftrightarrow\quad
 N-N'\in10\mathbb Z^{36}.
\]
In particular, the \(36\) sector residues determine exactly
the \(25\)-coordinate output. Preservation of quotients and
provenance constitutes additional information in the record.
\end{proposition}

\begin{proof}
Equality modulo \(10\) produces the same blocks and therefore
the same transverse coordinates. Conversely, if the transverse
coordinates agree, \(D_3(z-z')=D_4(z-z')=0\). Steps
\(3\) and \(4\) generate the cycle of length \(12\), so
\(z-z'\) is constant. Equality of total charges annihilates that
constant. Finally, the decimal representation of an integer between
\(0\) and \(999\) has three unique digits. Thus the
three residues of every window are recovered. The statement is an algebraic identity
for all \(N,N'\), not an inference obtained from finite tests.
\end{proof}

\subsection{Conjugation of counts and boundary terms}
\label{subsec:registro-conjugacion-incidencial}

Let \(\rho(m)=13-m\). Consider an incidence conjugation
that reverses the window order, preserves the arithmetic and
boundary classes of visits, and exchanges the trace channels.
Then
\[
 (A_m^\dagger,C_m^\dagger,V_m^\dagger)
 =(A_{\rho(m)},-C_{\rho(m)},V_{\rho(m)}).
\]
If \(c_m=[C_m]_{10}\), decimal reduction produces the identity
\begin{equation}
 z_m^\dagger=z_{\rho(m)}
       +100\mathbf1_{\{c_{\rho(m)}\ne0\}}-20c_{\rho(m)}.
 \label{eq:registro-conjugacion-decimal}
\end{equation}
Indeed, \([-C]_{10}=0\) when \([C]_{10}=0\), and
\([-C]_{10}=10-[C]_{10}\) otherwise. The formula preserves
exactly this difference. Negating the channel is not equivalent to negating
the entire decimal block.

Applying this identity to a concrete TPK involution requires
checking its action on incidences. For an occupation reader
evaluated before advancement, reversing a route
\(\gamma=(x_0,\ldots,x_n)\) satisfies
\[
 \sum_{k=0}^{n-1}f(x_{n-k})
 =\sum_{k=0}^{n-1}f(x_k)+f(x_n)-f(x_0).
\]
Endpoint terms are incorporated before reduction modulo
\(10\). On closed routes they vanish; in open windows they can
alter the residues. This correction, already present in the bidirectional
transport development, specifies what the connection
between route and counts must preserve.

\begin{proposition}[Necessity of the record's nonperiodic information]
Suppose the observable of a fixed route configuration satisfies
\(o(t+54)=o(t)\), multiplicity is preserved under that return,
and the same local reader acts in windows of \(9\) instants.
Then its counts and blocks satisfy
\[
 (A,C,V)_{m+6}=(A,C,V)_m,\qquad z_{m+6}=z_m.
\]
\end{proposition}

\begin{proof}
Translation \(t\mapsto t+54\) bijectively sends \(E_m\)
to \(E_{m+6}\), preserving each counted incidence and its multiplicity.
Equality passes to the sum and decimal reduction.
\end{proof}

Thus a presentation with \(12\) distinct blocks requires
the reader to preserve some information that does not factor through that
period-\(54\) observable: incidence selection, effective orientation,
endpoints, multiplicity, or memory. The proposition delimits an
inadmissible state reduction; it does not assert that the full dynamics has
period \(54\) or itself determine its terminal reader.


\subsection{Two readings of the terminal state}
\label{subsec:k-terminal}

Denote by \(x_{\mathrm{term}}\) the state produced by the composition
of selection, transport, and update up to the terminal cut. The
operator expression is
\[
 x_{\mathrm{term}}=
 (\mathcal U_{108}\circ\cdots\circ\mathcal U_1)(x_{\mathrm{seed}}),
 \qquad \mathcal U_t=\mathsf{Upd}_t\circ\mathsf{Tra}_t\circ\mathsf{Sel}_t.
\]
Its two relevant readings are
\[
 x_{\mathrm{term}}
 \xrightarrow{(\pi_{\mathrm{term}},R_{\mathrm{sgn}})}
 (B_{\mathrm{term}},U).
\]
The matrix \(B_{\mathrm{term}}\) retains the visible ternary reading; the
vector \(U\) preserves another coordinate of the same history. No arrow
\(B_{\mathrm{term}}\to U\) is inserted between them.

The terminal selector defined by the composition
\(\mathcal U_{108}\circ\cdots\circ\mathcal U_1\) and evaluated by the
incidence ledger of Subsection~\ref{subsec:registro-evaluacion-incidencial}
operates on catalogs of sizes 196, 197, and 196. The local signatures reduce their
\(196\cdot197\cdot196=7\,567\,952\) triples to 331, and the column
margin \(q_{\partial}=(1,2,2,1,2,0)\) leaves two matrices:
\[
 B_0=\begin{pmatrix}0&2&0&2&2&0\\0&0&1&2&2&0\\1&0&1&0&1&0\end{pmatrix},
 \qquad
 B_1=\begin{pmatrix}0&2&1&0&2&0\\0&0&1&2&2&0\\1&0&0&2&1&0\end{pmatrix}.
\]
The row orientation is \(\sigma=(1,1,-1)=(1,1,2)\) in
\(\mathbb F_3^3\). The row charges of the candidates are
\((0,2,0)\) and \((2,2,1)\). Only the second belongs to
\(\langle\sigma\rangle=\{(0,0,0),(1,1,2),(2,2,1)\}\); hence
\(B_{\mathrm{term}}=B_1\). This check proves the final step of the
selector from the declared census. It does not make the column margin
a consequence of the row charge, nor replace the preceding enumeration
with inspection of these two matrices.

The signed reader has the typed composition
\begin{equation}
 R_{\mathrm{sgn}}
 =\Pi_H\circ\mathcal E_{108}^{90,120}\circ\mathcal R_{12}:
 \mathcal X_{\mathrm{TPK}}^{\mathrm{term,mem}}
 \longrightarrow\mathbb Z^{12}.
 \label{eq:k-lector-firmado}
\end{equation}
The channel record determines
\[
 \mathcal E_{108}^{90,120}(\mathcal R_{12}(x_{\mathrm{term}}))
 =(b^{90},b^{120},Q_{\mathrm{TPK}}),
\]
with coordinates
\begin{align*}
 b^{90}={}&(-495,-116,-684,108,601,-90,\\
            &\hspace{19mm}-173,-88,313,560,-397,461),\\
 b^{120}={}&(-425,-281,-481,671,-255,30,\\
             &\hspace{19mm}475,-543,680,251,6,-128),\\
 Q_{\mathrm{TPK}}={}&6263.
\end{align*}
These quantities are used as coordinates of the preceding record, not
as parameters fitted through alpha. The reconstruction proved
in the following subsections starts exactly at this typed output.
The preceding evaluation is developed in
\S\ref{subsec:registro-evaluacion-incidencial}: visits and traces
determine the counts, which produce the blocks through
\eqref{eq:registro-bloque-incidencial} and the differences and charge through
\eqref{eq:registro-composicion-incidencial}. Proposition
\ref{prop:registro-36-residuos} exactly characterizes the
arithmetic information in that evaluation; subsection
\ref{subsec:registro-conjugacion-incidencial} preserves its boundary
terms under conjugation. Identification of the effective incidence collection
of this terminal state has the scope specified there.
The terminal domain, the incidence evaluation, the boundary conjugation,
and the arithmetic characterization producing the typed output are thus
identified within the article. The harmonic reconstruction below begins
with that output and preserves the preceding evaluation.

\subsection{From the channel record to the signed vector}
\label{subsec:k-pi-h}

The harmonic reader \(\Pi_H\) can be written without using either \(K\) or
alpha. To avoid confusing its orbital sums with the alpha blocks,
we denote them by \(s_1,s_2,s_3\). Let
\[
 B_r=\sum_{j=0}^{3}b^{120}_{r+3j},\qquad
 d_{r,j}=b^{90}_{r+3j},\qquad1\leq r\leq3,
\]
and
\begin{equation}
 s_1=\frac{Q_{\mathrm{TPK}}+2B_1+B_2}{3},\qquad
 s_2=s_1-B_1,\qquad s_3=s_2-B_2.
 \label{eq:k-sumas-armonicas}
\end{equation}
Division by three belongs to the three-orbit reader. Its
integrality is checked on compatible record outputs; it is not
assumed for an arbitrary element of \(\mathbb Z^{25}\).

The three signed blocks are
\begin{equation}
 U^{(r)}=(s_r,d_{r,1}-d_{r,3},d_{r,0}+d_{r,2},d_{r,0}-d_{r,2}).
 \label{eq:k-bloques-firmados}
\end{equation}
Integer substitution gives
\[
 (B_1,B_2,B_3)=(972,-1073,101),\qquad
 (s_1,s_2,s_3)=(2378,1406,2479),
\]
and therefore
\begin{align*}
 U^{(1)}&=(2378,-452,-668,-322),\\
 U^{(2)}&=(1406,998,-204,-28),\\
 U^{(3)}&=(2479,-551,-371,-997).
\end{align*}
Interleaving by columns yields
\begin{equation}
 U=(2378,1406,2479,-452,998,-551,
       -668,-204,-371,-322,-28,-997).
 \label{eq:k-u-firmado}
\end{equation}
The presence of negative values and components greater than 999
shows that \(U\) is not a base-thousand digit word. Nor is it
\(U_6\Vert U_6\), a catalog concatenation of period six. The difference
concerns domain and preserved information, not just values.

The constructive order so far is
\[
 x_{\mathrm{term}}\longrightarrow\mathcal R_{12}
 \longrightarrow(b^{90},b^{120},Q_{\mathrm{TPK}})
 \longrightarrow U.
\]
The inversion below produces the dodecaphasic record. The subsequent ability to recover
the channels from the dodecaphasic record is a reversibility check; it does not
reverse this causal order.

\subsection{Hadamard transformation by orbits and integer reconstruction}
\label{subsec:k-hadamard}

On the twelve-position cycle, step three determines three orbits:
\[
 (1,4,7,10),\qquad(2,5,8,11),\qquad(3,6,9,12).
\]
The separation into four sectors of each orbit is what, once
origin and orientation are fixed, admits the angular reading of ninety degrees. The
transformation does not act on three consecutive four-element sets in the natural order.

Let
\[
 H_4=\begin{pmatrix}
 1&1&1&1\\1&1&-1&-1\\1&-1&1&-1\\1&-1&-1&1
 \end{pmatrix}.
\]
If \(P_{\mathrm{orb}}\) reorders according to the three indicated orbits,
define
\[
 \mathcal H_{12}=P_{\mathrm{orb}}^{-1}
 (H_4\oplus H_4\oplus H_4)P_{\mathrm{orb}}.
\]

\begin{lemma}[Hadamard inversion and integrality]
\label{lem:k-hadamard}
We have
\[
 H_4^2=4I_4,\qquad H_4^{-1}=\frac14H_4,\qquad\det H_4=-16.
\]
For \(u\in\mathbb Z^4\), the preimage \(H_4^{-1}u\) is integral if and
only if \(H_4u\equiv0\pmod4\).
\end{lemma}

\begin{proof}
The rows are orthogonal and have squared norm four; the matrix is
symmetric. This gives \(H_4^2=4I_4\) and the inverse. Integer elimination,
or expansion of the determinant, gives the negative sign and value \(-16\).
The integrality condition is exactly divisibility by four
of all four coordinates of \(H_4u\).
\end{proof}

\begin{definition}[Integral dodecaphasic transformation]
The sublattice of signed records and its integral coordinatization are
\[
 \mathcal L_H=\{u\in\ZZ^{12}:\mathcal H_{12}u\equiv0\pmod4\},
 \qquad \mathscr K(u)=\tfrac14\mathcal H_{12}u.
\]
The map \(\mathscr K:\mathcal L_H\to\ZZ^{12}\) is an isomorphism
of abelian groups, with inverse \(k\mapsto\mathcal H_{12}k\).
The canonical record \(K\) is the image of the signed record generated
by the preceding operations, with phase origin and orientation fixed.
\end{definition}

\begin{proof}
The congruence defines exactly the integrality of \(\mathscr K(u)\).
The identity \(\mathcal H_{12}^2=4I\) proves the two inverse
compositions. In particular, every \(k\in\ZZ^{12}\) has preimage
\(\mathcal H_{12}k\in\mathcal L_H\).
\end{proof}

\begin{proposition}[Canonical integral record]
\label{prop:k-sello}
The inversion \(K^{(r)}=H_4U^{(r)}/4\) of the blocks in
\eqref{eq:k-bloques-firmados} produces
\begin{align*}
 K^{(1)}&=(234,729,621,794),\\
 K^{(2)}&=(543,659,58,146),\\
 K^{(3)}&=(140,824,914,601).
\end{align*}
In the natural order, the dodecaphasic record is
\begin{equation}
 K=(234,543,140,729,659,824,621,058,914,794,146,601).
 \label{eq:k-vector}
\end{equation}
It is the unique rational preimage of \(U\), and belongs to
\(\{0,\ldots,999\}^{12}\).
\end{proposition}

\begin{proof}
The three products before division are
\begin{align*}
 H_4U^{(1)}&=(936,2916,2484,3176),\\
 H_4U^{(2)}&=(2172,2636,232,584),\\
 H_4U^{(3)}&=(560,3296,3656,2404).
\end{align*}
All coordinates are divisible by four. Dividing and undoing the
orbit permutation gives \eqref{eq:k-vector}. Uniqueness follows from
rational invertibility, not from a search among candidates close to
a physical value.
\end{proof}

Integrality can also be formulated in lattice terms. The
determinantal divisors of \(H_4\) are \(1,2,4,16\): the entries
have greatest common divisor one; the order-two minors are even, with
at least one of absolute value two; the order-three minors are multiples of four, with
at least one of absolute value four. Hence
\[
 \operatorname{SNF}(H_4)=\operatorname{diag}(1,2,2,4),
\]
and for the global transformation,
\[
 \operatorname{coker}\mathcal H_{12}
 \cong(\mathbb Z/2\mathbb Z)^6\oplus(\mathbb Z/4\mathbb Z)^3.
\]
The index of the image is \(16^3=4096\). Thus not every signed
integer vector admits an integer preimage; the vector obtained satisfies the
precise congruences required by inversion. The factor \(1/4\) is
forced by the matrix and does not have the status of a free coefficient.

\subsection{Cyclic differences and the uniform component}
\label{subsec:k-transversal}

With cyclic indices modulo twelve, let
\[
 (D_3z)_m=z_m-z_{m+3},\qquad(D_4z)_m=z_m-z_{m+4},
 \qquad Q(z)=\sum_{m=1}^{12}z_m.
\]
The first operator compares sectors separated by three positions; the
second compares sectors separated by four. The channel names 90 and 120 refer here to
these separations in the oriented cycle, not to an identification with a
physical operator from another domain.

\begin{proposition}[Completeness of the transverse system]
\label{prop:k-transversal}
The operator
\[
 \mathcal D=(D_3,D_4,Q):\mathbb Z^{12}\longrightarrow\mathbb Z^{25}
\]
has rank twelve and nonzero invariant factors
\[
 \operatorname{diag}(1,\ldots,1,12),
\]
with eleven units. In particular, its outputs determine a unique dodecaphasic record.
\end{proposition}

\begin{proof}
If \(D_3z=D_4z=0\), the vector is constant along both steps.
Since \(3\) and \(4\) generate the cycle modulo twelve,
\(\ker D_3\cap\ker D_4=\langle\mathbf1\rangle\). Total charge
eliminates this last freedom because \(Q(\mathbf1)=12\).

For the integral claim, the rows of \((D_3,D_4)\) are oriented
incidences of a connected graph. A spanning tree provides eleven
unit invariant factors. Every order-twelve minor must contain the
row \(Q\). Adding the first eleven columns to the last makes the latter
zero in the incidence rows and twelve in the \(Q\) row. All
maximal minors are divisible by twelve, and the tree produces one of
absolute value exactly twelve. The last factor is therefore twelve.
\end{proof}

Checking \eqref{eq:k-vector} gives
\[
 D_3K=b^{90},\qquad D_4K=b^{120},\qquad Q(K)=6263.
\]
Formulas \eqref{eq:k-sumas-armonicas} and
\eqref{eq:k-bloques-firmados} in turn reconstruct \(U\) from these
differences. This closes a triangle of exact representations of the
same object: transverse record, signed vector, and dodecaphasic record. Equality
of results does not collapse their codomains into one; it proves that
the declared maps recover the relevant information.

% New formalization integrated in revision III; originals preserved.
% Source owners: U016, equations extractor-diferencias-k and lector-armonico-*;
% registro_k.tex, subsections k-pi-h, k-hadamard, and k-transversal.
\subsection{Integral compatibility and determination of the transverse record}
\label{img09:seccion}

The domain of this construction is the representation of the record of an
APP--TRIT--TPK history. Evaluation of record events,
compatibility of its channels, and inversion of its
coordinates are distinguished. The following result gives a complete integral criterion for
the second level and an explicit factorization connecting the first to it.
The proof uses the record operators and establishes their
relations for every element of the indicated integer domains.

\subsubsection{Characterization by adjacent increments}

The indices belong to $\mathbb Z/12\mathbb Z$. Let
\[
 (Sz)_i=z_{i+1},\qquad D_3=I-S^3,\qquad D_4=I-S^4,
 \qquad Q(z)=\sum_{i=0}^{11}z_i.
\]
Retain the transverse extractor
\[
 \mathcal D:\mathbb Z^{12}\longrightarrow\mathbb Z^{25},
 \qquad k\longmapsto(D_3k,D_4k,Q(k)).
\]
For a triple of integer coordinates $(b,c,q)$, write
\begin{equation}
 w_i=c_i-b_{i+1},\qquad
 P_0=0,\quad P_i=\sum_{j=0}^{i-1}w_j\ (1\leq i\leq12),
 \qquad T(w)=\sum_{i=0}^{11}P_i.
 \label{img09:incremento}
\end{equation}

\begin{theorem}[Integral image of the transverse extractor]
\label{img09:imagen}
A triple $(b,c,q)\in\mathbb Z^{25}$ belongs to
$\operatorname{im}\mathcal D$ if and only if it satisfies
\begin{align}
 b_i&=w_i+w_{i+1}+w_{i+2}&& (0\leq i<12),
 \label{img09:relaciones}\\
 \sum_{i=0}^{11}w_i&=0,
 \label{img09:cierre}\\
 q+T(w)&\equiv0\pmod {12}.
 \label{img09:congruencia}
\end{align}
Its unique integer preimage is
\begin{equation}
 k_i=\frac{q+T(w)}{12}-P_i,\qquad0\leq i<12.
 \label{img09:inversa}
\end{equation}
The twelve equations of \eqref{img09:relaciones} and equation
\eqref{img09:cierre} are linearly independent over $\mathbb Q$.
\end{theorem}
\begin{proof}
For an output of $\mathcal D$, one has
\[
 c_i-b_{i+1}=(k_i-k_{i+4})-(k_{i+1}-k_{i+4})=k_i-k_{i+1}.
\]
The sum of three increments gives $b_i$; their full cyclic sum
is zero. Moreover, $P_i=k_0-k_i$, so
$q=12k_0-T(w)$. All conditions and the inverse formula follow.

Conversely, \eqref{img09:congruencia} makes formula
\eqref{img09:inversa} integral; \eqref{img09:cierre} allows its cyclic
prolongation. Its adjacent differences are $w_i$. Thus,
$D_3k=b$ and
\[
 (D_4k)_i=w_i+w_{i+1}+w_{i+2}+w_{i+3}
 =w_i+b_{i+1}=c_i.
\]
Summing the inverse gives $Q(k)=q$. The same formula proves uniqueness.
For independence, the invertible integer change
$(b,c)\mapsto(b,w=c-Sb)$ turns the first twelve relations into
$b=(I+S+S^2)w$. Each solves for a different coordinate of $b$.
The condition $\sum w_i=0$ is an additional independent relation.
\end{proof}

The rational image has dimension twelve. Its integral structure adds
a congruence of order twelve: this is the constructive manifestation of
the final Smith factor of $\mathcal D$. The eleven increments
$w_0,\ldots,w_{10}$ and the charge $q$ are sufficient coordinates, with
$w_{11}=-\sum_{i=0}^{10}w_i$ and congruence
\eqref{img09:congruencia}. The original twenty-five coordinates
preserve these same twelve degrees of freedom through thirteen relations.

\subsubsection{Commutativity with the harmonic reading}

Let $\mathcal H_{12}$ be the record's Hadamard transformation,
with orbits $(r,r+3,r+6,r+9)$ for $r=0,1,2$. Its block is
\[
 H_4=\begin{pmatrix}
 1&1&1&1\\1&1&-1&-1\\1&-1&1&-1\\1&-1&-1&1
 \end{pmatrix},\qquad H_4^2=4I_4.
\]
For compatible $(b,c,q)$, define
\[
 B_r=\sum_{j=0}^3c_{r+3j},\qquad
 A_0=\frac{q+2B_0+B_1}{3},\quad A_1=A_0-B_0,
 \quad A_2=A_1-B_1.
\]
The reader $\Pi_H$ of \eqref{eq:k-bloques-firmados} has blocks
\[
 (\Pi_H(b,c,q))^{(r)}=
 (A_r,b_{r+3}-b_{r+9},b_r+b_{r+6},b_r-b_{r+6}).
\]

\begin{proposition}[Agreement of the two reconstructions]
\label{img09:triangulo}
On $\operatorname{im}\mathcal D$,
\[
 \Pi_H\mathcal D=\mathcal H_{12},\qquad
 \mathcal D^{-1}=\tfrac14\mathcal H_{12}\Pi_H.
\]
The image of $\Pi_H$ restricted to that domain is exactly
\[
 \mathcal L_H=
 \{u\in\mathbb Z^{12}:\mathcal H_{12}u\equiv0\pmod4\}.
\]
\end{proposition}
\begin{proof}
For $k=\mathcal D^{-1}(b,c,q)$, let
$a_r=\sum_{j=0}^3k_{r+3j}$. Shifting four positions
sends each orbit to the next, so $B_r=a_r-a_{r+1}$.
Together with $q=a_0+a_1+a_2$, these identities give $A_r=a_r$.
On an orbit, write $(x_0,x_1,x_2,x_3)$ for the coordinates
of $k$ and $d_j=x_j-x_{j+1}$. Then
\[
 d_1-d_3=x_0+x_1-x_2-x_3,\quad
 d_0+d_2=x_0-x_1+x_2-x_3,\quad
 d_0-d_2=x_0-x_1-x_2+x_3.
\]
These are the remaining three rows of $H_4k^{(r)}$. Inversion and the
integral characterization follow from $\mathcal H_{12}^2=4I$.
\end{proof}

The condition $\Pi_H(b,c,q)\in\mathcal L_H$ controls only the
harmonic output. Membership of the twenty-five coordinates in
$\operatorname{im}\mathcal D$ additionally requires
\eqref{img09:relaciones}--\eqref{img09:congruencia}.
For example, adding the vector $e_0-e_3$ to $c$ preserves the three
orbital sums and leaves $\Pi_H$ unchanged. Starting from $(0,0,0)$,
one thus obtains a triple with $\Pi_H=0$ that violates
\eqref{img09:relaciones}. This is an exact negative control of
channel compatibility, independent of any terminal evaluation.

\subsubsection{Determination and transport on the record's orbits}

Theorem~\ref{img09:imagen} determines $k$ once its
channels and charge have been produced. As a condition on outputs not yet
individualized, the same image is a group isomorphic to
$\mathbb Z^{12}$. Fixing only a charge $q$, the solutions
form an affine class of
\[
 L_0=\{v\in\mathbb Z^{12}:\sum v_i=0\}
 =\bigoplus_{i=0}^{10}\mathbb Z(e_i-e_{11}).
\]
For example, $100\mathbf1+e_0$ and $100\mathbf1+e_1$ are two
different records of charge $1201$, with components in
$\{0,\ldots,999\}$. Each satisfies every image condition
and every Hadamard congruence through its own
channels. These witnesses compare compatibility equations;
they are not attributed the status of terminal HMT histories.

\begin{proposition}[Cyclic covariance and stabilizer of the record]
\label{img09:covariancia}
Let $\rho$ be window translation on a record domain
$\mathscr R$, and $S$ the preceding translation on $\mathbb Z^{12}$.
On the orbit of $R\in\mathscr R$, a covariant reader is
determined by a vector
\[
 k\in\operatorname{Fix}(S^p),
\]
where $p\mid12$ is the least period of $R$ under $\rho$.
The charge condition $Q(k)=q$ admits integer solutions if and only if
$12/p$ divides $q$; when they exist, they form an affine class of rank
$p-1$. In particular, on a free orbit, covariance and charge
leave eleven degrees of freedom.
\end{proposition}
\begin{proof}
The rule $F(\rho^jR)=S^jk$ is well defined exactly when
$S^pk=k$. Such vectors repeat a word of length
$p$ exactly $12/p$ times. Their charge is $\frac{12}{p}\sum_{i=0}^{p-1}k_i$.
Divisibility and rank follow from this formula. Channel
covariance follows from $D_aS=SD_a$ and $QS=Q$.
On $\mathcal L_H$, the corresponding action is
$\widehat S_H=\mathcal H_{12}S\mathcal H_{12}/4$;
by Proposition~\ref{img09:triangulo}, both reconstructions
intertwine these actions.
\end{proof}

If the domain carries a marked origin, the action transporting
that mark as well must be used. The proposition concerns this declared
cyclic action; it does not assume periodicity of the enriched
history. Naturality under truncations after the
$108$ cutoff is obtained, for any already defined reader $F_{108}$, through
$F_N=F_{108}\circ\tau_{N,108}$. This naturality preserves the initial
evaluation rule; it does not by itself choose one of the possible rules
on the first one hundred eight events.

\subsubsection{Event evaluation and composition of contributions}

The preceding characterization provides a factorization of extraction through two
effective readings of the record:
\begin{equation}
 \omega:\mathscr R\longrightarrow
 \{w\in\mathbb Z^{12}:\sum w_i=0\},\qquad
 q_{\rm reg}:\mathscr R\longrightarrow\mathbb Z,
 \quad q_{\rm reg}(R)+T(\omega(R))\equiv0\pmod {12}.
 \label{img09:contrato-productor}
\end{equation}
These maps must be evaluated on the events, orientation,
incidence, and memory produced by APP--TRIT--TPK. Once their
actions on these data are given, the composition
\[
 R\longmapsto(w,q)\longmapsto
 ((I+S+S^2)w,(I+S+S^2+S^3)w,q)
 \xrightarrow{\Pi_H}u\xrightarrow{\mathcal H_{12}/4}k
\]
is determined and satisfies all inverse checks. The formula
$w=c-Sb$ also allows connection to a producer already delivering both
channels: it is a check on its output, prior to any comparison
with the terminal witness.

\begin{proposition}[Accumulation of contributions and prospective uniqueness]
\label{img09:acumulacion}
Suppose the initial state and each event $a_t$ of the record
produce, through previously fixed rules, contributions
$(\delta w_t,\delta q_t)$ with
\[
 \sum_i\delta w_{t,i}=0,\qquad
 \delta q_t+T(\delta w_t)\equiv0\pmod {12}.
\]
With compatible initial conditions $(w^{(0)},q^{(0)})$, the
update
\[
 w^{(t+1)}=w^{(t)}+\delta w_t,\qquad
 q^{(t+1)}=q^{(t)}+\delta q_t
\]
produces a unique record at each stage. Its increment is
\[
 \delta k_{t,i}=
 \frac{\delta q_t+T(\delta w_t)}{12}-P_i(\delta w_t).
\]
The construction respects concatenation of records, and the reading
of a prefix remains unchanged under subsequent prolongations.
\end{proposition}
\begin{proof}
The compatibility conditions are preserved under addition. Formula
\eqref{img09:inversa} is additive in $(w,q)$ and produces the
indicated increment. Induction fixes each state successively.
The sum over two segments is the concatenated sum, and a prefix receives
exactly the same summands before and after prolonging the history.
\end{proof}

The proposition provides a sufficient criterion, not a requirement of
linearity for every TPK reader. If events are transported by nontrivial
actions, their contributions can be expressed in the terminal chart
through the corresponding affine cocycle. The cocycle equation then
preserves the same independence from the partition of the trajectory.
In both cases, event rules and the initial state constitute the
constructive data that must come from the ledger: choosing them to
reproduce an expected vector would replace that provenance with another task.

The result of this development is the exact determination condition
\eqref{img09:contrato-productor}, its proved inversion, and its transport
identities. Its scope is the transverse record; the remaining
HMT constructions retain their own domains and proofs.


\subsection{Two rational readings: \texorpdfstring{\(\kappa_{36}\)}{kappa36} and \texorpdfstring{\(\kappa_{\mathrm{per}}\)}{periodic kappa}}
\label{subsec:k-kappas}

Once \(K\) has been constructed, the base-thousand chart permits two
distinct operations. The first reads a single turn; the second prolongs the dodecaphasic record
periodically. Let \(B=1000\), and define its concatenation integer
\begin{equation}
 N_K=\sum_{m=1}^{12}K_mB^{12-m}
 =234543140729659824621058914794146601.
 \label{eq:k-numerador}
\end{equation}
The leading zeros of a block, such as that of 058, are part of its
three-digit representation. They do not change its integer value, but allow
the twelve blocks to be concatenated unambiguously.

\begin{definition}[Finite and periodic readings of the dodecaphasic record]
\label{def:k-kappas}
The one-turn reading is
\[
 \kappa_{36}=\sum_{m=1}^{12}K_mB^{-m}=\frac{N_K}{B^{12}}.
\]
The periodic prolongation is
\[
 K_j^{\mathrm{per}}=K_{1+((j-1)\bmod12)},\qquad
 \kappa_{\mathrm{per}}=\sum_{j\geq1}K_j^{\mathrm{per}}B^{-j}.
\]
\end{definition}

\begin{proposition}[Exact value and period of the prolonged dodecaphasic record]
\label{prop:k-periodo}
We have
\begin{equation}
 \kappa_{\mathrm{per}}=\frac{N_K}{B^{12}-1},\qquad
 \kappa_{\mathrm{per}}-\kappa_{36}
 =\frac{N_K}{B^{12}(B^{12}-1)}
 =B^{-12}\kappa_{\mathrm{per}}.
 \label{eq:k-diferencia-kappas}
\end{equation}
The infinite word of the dodecaphasic record has least period twelve in base thousand and
least period thirty-six in base ten.
\end{proposition}

\begin{proof}
The sum of the first turn is \(N_KB^{-12}\); each new turn
multiplies its contribution by \(B^{-12}\). The geometric series yields
\(N_K/(B^{12}-1)\), and subtraction gives the second identity.

The twelve components of \(K\) are distinct. A proper period of the cycle
would identify two of them; the least block period is therefore
twelve. The decimal expansion is canonical, since the dodecaphasic record has no tail
of 999 blocks. If its least decimal period were \(d\), it would divide
36. Grouping digits in threes would produce a block period
\(d/\gcd(d,3)\). No proper divisor of 36 gives twelve through this
operation. Hence the least decimal period is 36.
\end{proof}

Both readings are rational, but they are not the same rational number. The first
terminates after twelve blocks; the second repeats the turn. Moreover,
\(0<\kappa_{\mathrm{per}}-\kappa_{36}<B^{-12}\). Their closeness does not
permit substituting one for the other in an exact identity. In particular,
the finite alpha closure uses \(\kappa_{36}\), whereas the
sum of the prolonged section uses \(\kappa_{\mathrm{per}}\).

\subsection{Reversibility of the rational constant and phase change}

The rational constant \(\kappa_{\mathrm{per}}\) determines
the entire record \(K\) when the base, length, and reading origin
are retained. Its numerical and structural functions are related
by an exact inversion.

\begin{proposition}[Integral recovery and cyclic transformation]
Let \(I(K)=\sum_{j=1}^{12}K_j1000^{12-j}\). Then
\[
 I(K)=(1000^{12}-1)\kappa_{\mathrm{per}}.
\]
The integer expansion of \(I(K)\) into twelve base-\(1000\) blocks
recovers every \(K_j\), including the leading zeros in each block.
If \(\rho(K)=(K_2,\ldots,K_{12},K_1)\), then
\begin{equation}
 \kappa(\rho K)=1000\kappa(K)-K_1.
 \label{eq:k-rotacion-racional}
\end{equation}
\end{proposition}
\begin{proof}
The first identity inverts the definition of periodic evaluation.
Iterated Euclidean division by \(1000\), with twelve fixed positions,
recovers the vector. Furthermore,
\[
 I(\rho K)=1000I(K)-K_1(1000^{12}-1).
\]
Dividing by \(1000^{12}-1\) yields \eqref{eq:k-rotacion-racional}.
\end{proof}

Recovery therefore composes \(\kappa\mapsto K\mapsto U\)
and the already constructed transformations \(D_3K,D_4K,Q\).
Rational evaluation preserves the signed record and the cyclic differences
that depend on it. The total depth and transport operators of
a history remain in their own coordinates: the domain of the preceding
inversion is the integral dodecaphasic record.

Write \(\kappa_j=\kappa(\rho^jK)\), with indices modulo \(12\).
The same identity gives
\[
 K_{j+1}=1000\kappa_j-\kappa_{j+1},\qquad
 \sum_{j=0}^{11}\kappa_j=\frac{\sum_{j=1}^{12}K_j}{999}
 =\frac{6263}{999}.
\]
Phase acts exactly on the constant: its value is fixed
in the canonical reading, while its rotations follow an affine law.
The orbit sum is invariant under the choice of origin.

The subsequent incidence also has an expression in these
coordinates. The cubic threshold determines
\[
 B_K=\{j+1:1000\kappa_j-\kappa_{j+1}\ge729\}
 =\{4,6,9,10\}.
\]
The tetrad is recovered from the rational orbit because the latter
first recovers the integer components. The signed support associated
with \(\alpha\) additionally uses its regional records and boundary
normalization. The relationship between coordinate and incidence is thus
expressed through identifiable operations.

This is an operational definition of the function of \(K\): it provides
a reversible integral coordinatization of the oriented record, whose
periodic evaluation and incidence functions participate in
subsequent constructions. The value, transformation, and domain
are preserved together.



\subsection{Periodicity and cokernel of the transport operator}
\label{subsec:k-clase-acarreo}

The periodicity just proved belongs to the dodecaphasic record. It does not imply
periodicity of the regional coordinates, the closure carries,
or alpha. It is useful to clarify this separation through a second,
purely integral quotient.

Let \(S\) be the cyclic shift of twelve coordinates, with fixed
orientation, and let \(b\geq2\) be an integer base. The periodic
carry operator is \(\partial_b=S-bI\). If a periodic identity is written
\[
 A=P+E-\Phi-K+\partial_bC,
\]
then the transformations
\[
 K\longmapsto K+\partial_bh,\qquad C\longmapsto C+h
\]
preserve its right-hand side. This invariance compares representatives of a
carry class; it does not authorize changing the dodecaphasic record already
selected by the record.

\begin{proposition}[Quotient of periodic carry]
\label{prop:k-cociente-acarreo}
With the preceding orientation,
\[
 \operatorname{coker}(S-bI)\cong\mathbb Z/(b^{12}-1)\mathbb Z.
\]
\end{proposition}

\begin{proof}
The cyclic module is presented as
\(\mathbb Z[t]/(t^{12}-1)\). Imposing \(S=bI\) is equivalent to imposing
\(t=b\) and leaves the relation \(b^{12}-1=0\). Evaluating the
coefficients at \(b\), with the chosen concatenation order, represents
the resulting class.
\end{proof}

The denominator \(B^{12}-1\) of \(\kappa_{\mathrm{per}}\) and this
quotient arise from the same cyclic length, but their objects remain
different: one is a rational evaluation and the other a quotient
of a carry module. The finite alpha closure will be an oriented
chain with a right boundary. Its prolongation will transport a
compatible remote boundary, not a carry identified by decree
between the first and thirteenth positions.

The next chapter is thus prepared. The dodecaphasic record has been fixed before
solving the alpha recurrence; its Hadamard reconstruction is
integral and unique; its transverse differences recover the record; and
its two rational evaluations have separate domains and uses. None
of these facts uses alpha as an input, and none yet determines the
arithmetic type of the output of the complete closure.
